\documentclass[11pt,a4paper]{article}
\usepackage[T1]{fontenc}
\usepackage{lmodern,amsmath,amssymb,amsthm,mathtools}
\usepackage[margin=27mm]{geometry}
\usepackage[colorlinks=true,linkcolor=blue,urlcolor=blue]{hyperref}
\newtheorem{theorem}{Theorem}
\newtheorem{lemma}[theorem]{Lemma}
\newtheorem{corollary}[theorem]{Corollary}
\newcommand{\Q}{\mathbb Q}\newcommand{\Z}{\mathbb Z}
\title{Scalar contact normalization by finite-field avoidance}
\author{}\date{8 November 2026}
\begin{document}\maketitle
\begin{abstract}
We construct a scalar contact of numerator degree twenty and
denominator degree forty that preserves the exact arithmetic
conditions of the specified rational source. A residue avoidance
condition ensures a uniform odd valuation at a fixed prime.
The required prime exponent is $137731$ in the numerator. The
construction uses the existing coordinates and finite algebra
generators. A general normalization lemma explains the method,
and exact finite computations certify all fixed residue claims.
\end{abstract}

\section{A general normalization lemma}

\begin{lemma}\label{lem:general}
Let $f\in\Z[X]$ have even degree $d>0$, and let $p$ be an odd
prime not dividing its leading coefficient. Suppose $f$ has no
zero modulo $p$, and that $f(x)\ge m>0$ for every real $x$.
Choose $a\in\mathbb F_p^\times$ such that
$-a^{-1}\notin f(\mathbb F_p)$. Let $C$ be a positive integer
congruent to $a$ modulo $p$, let $L$ be a positive odd integer,
and let $c\in\Z$ be nonzero and prime to $p$. Assume
$C^2m>c^2p^L$. Then
\[
 R_f(X)=\frac{c^2p^L f(X)}{(1+C f(X))^2}
\]
has the following properties for every rational $X$:
\[
 0<R_f(X)<1,\qquad
 v_p(R_f(X))=
 \begin{cases}
 L,&v_p(X)\ge0,\\
 L-dv_p(X),&v_p(X)<0.
 \end{cases}
\]
In particular its $p$-valuation is odd and at least $L$.
Its numerator and denominator, in lowest terms over $\Q[X]$,
have degrees $d$ and $2d$.
\end{lemma}
\begin{proof}
The real denominator is larger than $C^2f(X)^2$, so
\[
 0<R_f<\frac{c^2p^L}{C^2f(X)}
       \le\frac{c^2p^L}{C^2m}<1.
\]
For a $p$-integral input, both $f(X)$ and $1+Cf(X)$ are
$p$-adic units by their residue assumptions. For an input of
valuation $n<0$, the leading term of $f(X)$ uniquely dominates
the other terms, giving $v_p(f(X))=dn$. Since $C$ is a unit,
$v_p(1+Cf(X))=dn$. The two formulas follow, and their parity
follows from the even degree. Finally
$\gcd(f,1+Cf)=1$ in $\Q[X]$.
\end{proof}

The missing residue $-a^{-1}$ is often available without any
arithmetic conjecture. For example, if $f(X)=g(X^e)$ modulo
$p$ and $\gcd(e,p-1)=q>1$, then
$|f(\mathbb F_p)|\le1+(p-1)/q$. If this upper bound is smaller
than $p-1$ and $f$ has no zero, there is a missing nonzero
residue. The coefficient $C$ can then be made arbitrarily
large in the required residue class; additional divisibility
conditions coprime to $p$ can be imposed by the Chinese
remainder theorem. Thus real normalization and fixed-prime
avoidance can be chosen independently.

\section{The explicit scalar contact}

Define
\[
 \beta(X)=216X^4+450459,
 \qquad r(X)=\beta(\beta-1)(\beta-2401)(32\beta+343)(2\beta-1).
\]
Put
\begin{equation}\label{eq:constants}
 \ell=76249,\quad L=137731,\quad
 C=17\bigl(676\ell^{68862}+1\bigr),\quad
 R(X)=\frac{102^2\ell^L r(X)}{(1+C r(X))^2}.
\end{equation}

\begin{lemma}\label{lem:residues}
The integer $\ell$ is prime. Both $r(X)$ and $1+17r(X)$ have
no zero in $\mathbb F_\ell$.
\end{lemma}
\begin{proof}
Trial division through $\lfloor\sqrt\ell\rfloor=276$ proves
primality. For a zero of $r$, one of the five values of $\beta$
would be $0,1,2401,-343/32,1/2$. The corresponding values of
$X^4=(\beta-450459)/216$ modulo $\ell$ are
\[
 32862,\quad32509,\quad24048,\quad904,\quad70810.
\]
Their powers with exponent $(\ell-1)/4$ are respectively
$76248,76248,76248,76248,40177$. None is zero or one, so none
is a fourth power.

For the second assertion the exact certificate supplied with
this paper checks
\[
 \gcd\bigl(1+17r(X),\ X^{76249}-X\bigr)=1
       \quad\hbox{in }\mathbb F_{76249}[X].
\]
It records the coefficients of the degree-twenty first
polynomial, the remainder of $X^{76249}-X$ modulo it, and a
B\'ezout identity with constant term one. Modular repeated
squaring verifies the remainder, and multiplication verifies
the B\'ezout identity. Since $X^{76249}-X$ is the product of
the distinct linear polynomials over this field, this proves
the assertion. An independent complete residue enumeration
gives the same result: $X^4$ has $19063$ values, and $r(X)$
has $16859$ values, none equal to $-17^{-1}$.
\end{proof}

\begin{lemma}\label{lem:small-primes}
For every integer $X$, $r(X)$ is a square in each of
$\Q_2,\Q_3,\Q_{17}$, with
\[
 v_2(r)=v_3(r)=2,\qquad v_{17}(r)\in\{0,2\}.
\]
\end{lemma}
\begin{proof}
One has $\beta\equiv3\pmod8$ and $\beta\equiv18\pmod{27}$.
The two even factors $\beta-1,\beta-2401$ have valuation one,
and their odd parts differ by $1200\equiv0\pmod8$. Their
odd-part product is therefore one modulo eight. The three
other factors have product $3\cdot7\cdot5\equiv1\pmod8$.
Thus $v_2(r)=2$ and $r/4\equiv1\pmod8$.
At three only $\beta$ is divisible by three, with valuation
two, and
\[
 r/9\equiv2\cdot2\cdot2\cdot1\cdot2\equiv1\pmod3.
\]
At seventeen the five possible fourth powers and their images
are
\[
 \begin{array}{c|rrrrr}
 X^4&0&1&4&13&16\\\hline
 \beta&10&5&7&13&15\\
 r(X)&0&15&2&13&2.
 \end{array}
\]
The nonzero entries in the last row are squares. In the first
column $17\mid X$ and
\[
 32\beta+343=6912X^4+17^2\cdot49879,
 \qquad49879\equiv1\pmod{17}.
\]
This factor has valuation two, and the other four factors have
unit product $10\cdot9\cdot6\cdot2\equiv9\pmod{17}$.
The usual local square criteria now give the claim.
\end{proof}

\begin{theorem}\label{thm:explicit}
The rational function in~\eqref{eq:constants} has the following
properties.
\begin{enumerate}
\item Its numerator and denominator in lowest polynomial terms
have degrees twenty and forty, respectively.
\item For all $X\in\Q$, $0<R(X)<1$ and
\[
 v_\ell(R(X))=
 \begin{cases}
 137731,&v_\ell(X)\ge0,\\
 137731-20v_\ell(X),&v_\ell(X)<0.
 \end{cases}
\]
\item If $p>76176$, $p\ne\ell$, and $v_p(r(X))>0$, then
$v_p(R(X))=v_p(r(X))$.
\item If $X\in\Z$, the reduced denominator of $R(X)$ is a
square, and $R(X)$ is a square in $\Q_2,\Q_3,\Q_{17}$, with
\[
 v_2(R)=v_3(R)=4,\qquad v_{17}(R)\in\{2,4\}.
\]
\end{enumerate}
\end{theorem}
\begin{proof}
Since $\beta\ge450459$ on $\mathbb R$, all five factors of
$r$ are positive and increasing as functions of $\beta$.
Consequently
\[
 r(X)\ge r(0)=1180712527837405404989927078052.
\]
The exact integer inequality
\[
 36\ell^7<676^2r(0)
\]
implies
$C^2r(0)>17^2\cdot676^2\ell^{137724}r(0)
>102^2\ell^{137731}$.
Now apply Lemma~\ref{lem:general} with $f=r$, $d=20$,
$a=17$, and $c=102$, using Lemma~\ref{lem:residues}.
This gives parts~(1) and~(2).

For part~(3), $C$ is an integer, so $v_p(Cr)>0$ and the
denominator is a unit. The multiplier $102^2\ell^L$ is a
$p$-adic unit for all primes in the stated range.

For integer input let $D_0=1+Cr(X)$. It is coprime to $r(X)$.
It is a unit at two and three because $r$ is divisible by four
and nine, a unit at seventeen because $17\mid C$, and a unit
at $\ell$ by Lemma~\ref{lem:residues}. The only prime factors
of $102^2\ell^L$ are these four primes. Therefore the displayed
numerator is coprime to $D_0$, so the reduced denominator is
exactly $D_0^2$. Finally
$\ell\equiv1\pmod8$, $\ell\equiv1\pmod3$, and
$\ell\equiv4\pmod{17}$, making $\ell$ a square in the three
completions. Lemma~\ref{lem:small-primes} and the square
multiplier $102^2$ give the stated squareclasses and depths.
\end{proof}

\section{Integration into the rational source}

The scalar extraction argument of the 29-coordinate source
requires a positive rational contact below one, an odd
$\ell$-valuation at least $137731$, preservation of consumed
positive contact valuations away from $\ell$, and, for integer
inputs, a square denominator and the preparation conditions at
two, three and seventeen. Theorem~\ref{thm:explicit} supplies
these same hypotheses. The finite support can change; the
completeness theorem permits preparation at any fixed finite
support after fixing the input. Consequently the replacement
preserves the source's soundness and completeness under exactly
its existing arithmetic premises. No coordinate or finite
algebra generator is added.

The final polynomial's degree and coefficient bounds must be
obtained from its revised expression graph. The degree-forty
denominator and the smaller fixed prime exponent are exact
improvements to the scalar source; they do not alone certify
a numerical global bound for the eliminated polynomial.

\clearpage
\section{A compact residue certificate and reproduction}

The root-avoidance calculation also has a small certificate in
the variable $U=X^4$. Define four polynomials over
$\mathbb F_{76249}$ by their coefficient vectors, in ascending
order of powers of $U$:
\[
\begin{array}{c|l}
 F_0&(2124,22608,14921,12900,56669,30599)\\
 G_0&(69700,50246,7543,69483,32972)\\
 A_0&(32273,39191,5614,60392,71693)\\
 B_0&(21082,70119,44908,68100,13136).
\end{array}
\]
Expansion gives $1+17r(X)=F_0(X^4)$. Repeated squaring and
polynomial multiplication give the exact identities
\[
 U^{19062}-1\equiv G_0(U)\pmod{F_0(U)},\qquad
 A_0F_0+UB_0G_0=1.
\]
If $F_0(x^4)=0$ for $x\in\mathbb F_{76249}$, then $x\ne0$
since $F_0(0)=2124$. The element $u=x^4$ satisfies
$u^{19062}=1$, so the first identity gives $G_0(u)=0$.
The second identity then gives $0=1$. This supplies the
root-avoidance argument with only degree-five arithmetic and
one modular exponentiation. The direct residue enumeration
is a separate independent verification.

The checker is
\begin{center}
\texttt{python check\_scalar\_normalization.py}
\end{center}
It uses integer arithmetic and polynomial arithmetic over the
prime field, with no floating-point tests. Its output file is
\begin{center}
\texttt{scalar-normalization-certificate.json}.
\end{center}
This includes the complete modular B\'ezout certificate. The direct enumeration
is independent of the modular polynomial Euclidean algorithm.
The valuation formulas for all rational inputs are proved in
Lemma~\ref{lem:general}; no bounded scan is substituted for
that argument.
\end{document}
